%%%%%%%%%%%%%%%%%%%%%%%%%%%%%%%%%%%%%%%%%
% Important note:
% This template requires the resume.cls file to be in the same directory as the
% .tex file. The resume.cls file provides the resume style used for structuring the
% document.
%%%%%%%%%%%%%%%%%%%%%%%%%%%%%%%%%%%%%%%%%

%----------------------------------------------------------------------------------------
%	PACKAGES AND OTHER DOCUMENT CONFIGURATIONS
%----------------------------------------------------------------------------------------

\documentclass{resume} % Use the custom resume.cls style

\usepackage[left=0.75in,top=0.6in,right=0.75in,bottom=0.6in]{geometry} % Document margins
\newcommand{\tab}[1]{\hspace{.2667\textwidth}\rlap{#1}}
\newcommand{\itab}[1]{\hspace{0em}\rlap{#1}}
\name{Felix Kimmerle} % Your name
\address{Tübinger Straße 10 \\ 70178 Stuttgart} % Your address
\address{(+49) 15737879503 \\ FelixKimmerle@hotmail.com} % Your phone number and email

\begin{document}

%----------------------------------------------------------------------------------------
%	EDUCATION SECTION
%----------------------------------------------------------------------------------------

\begin{rSection}{Ausbildung}
{\bf Universität Stuttgart} \hfill {\em 14.10.2019 - voraussichtlich November 2024} 
\\ Bachelor of Science Informatik \hfill {\em aktueller Schnitt: 1,9}
\\ Anwendungsfach: Maschinelle Sprachverarbeitung (Parsing, Information Retrieval and Text Mining)
\\ Auswahl relevanter Module: Systemkonzepte und -programmierung, Programmierparadigmen, Computergraphik, Numerische Grundlagen


{\bf Oscar-Paret-Schule Freiberg am Neckar} \hfill {\em 2011 - 2019} 
\\ Allgemeine Hochschulreife \hfill {\em Note: 2,5}
\\ Leistungskurse: Mathematik, Physik, Biologie, Deutsch
\end{rSection}

%----------------------------------------------------------------------------------------
%	EXPERIENCE SECTION
%----------------------------------------------------------------------------------------

\begin{rSection}{Berufserfahrung}
{\bf Bachelorarbeit, Robert Bosch GmbH}{, VM-VDS/EDS6, IPVS} \hfill {\em 15.05.2024 - voraussichtlich November 2024}\\
Entwicklung von Vor- und Nachverarbeitungsmethoden zur Simulation induktiver Positionssensoren sowie Erweiterung der BEM-BEM-Simulation zur Unterstützung mehrerer Subdomains und Topologien mit Löchern.

{\bf Werkstudent, Robert Bosch GmbH}{, VM-VDS/EDS6} \hfill {\em 14.10.2019 - 15.05.2024}\\
Programmierung einer 3D-Raytracing-Simulation in C++; Weiterentwicklung eines Datenbankvisualisierers in Python mit Qt; Optimierung und Parallelisierung von Simulationsauswertungen; Entwicklung von Tools zur Computerverwaltung; Erweiterung der Python-Bibliothek meshio, um die Unterstützung von Timesteps in Exodus II-Dateien zu ermöglichen, wodurch die Auswertung von Edyson2020-Simulationsergebnissen vereinfacht wurde; Einrichtung und Betreuung eines Linux-Clusters mit SLURM.

{\bf Tutor, Universität Stuttgart}{, SQA, SEC} \hfill {\em 01.05.2020 - 19.08.2024}
\begin{itemize}
    \item Tutor für Informatik II (Programmierung in C++): Betreuung von Übungsgruppen, Klausurvorbereitung, Erstellung von Aufgaben.
    \item Tutor für Programmierung und Softwareentwicklung: Korrektur und Besprechung von Hausaufgaben, Erstellung von Aufgaben.
\end{itemize}

{\bf Hiwi, Universität Stuttgart}{, TIK} \hfill {\em 01.08.2022 - 01.05.2024}\\
Beratung zu TIK-Diensten, Beantwortung und Delegierung von Support-Tickets, Aufsicht im Makerspace.
\end{rSection}

%----------------------------------------------------------------------------------------
%	SKILLS SECTION
%----------------------------------------------------------------------------------------

\begin{rSection}{Fähigkeiten}
{\bf Programmiersprachen}: C, C++ (seit 2014), C\#, Java, Python, Lua, RISC-V, bash \\
{\bf Betriebssysteme und Tools}: Linux (LPIC-Essentials Zertifikat), Windows, WPF, QT, GCC, Clang, MSVC, Makefiles, CMake, Visual Studio Code, Visual Studio, IntelliJ IDEA, Eclipse, OpenGL, SFML, SDL, WSL, QEMU, SLURM, Vivado \\
{\bf Sprachen}: Deutsch (Muttersprache), Englisch (B2), Russisch (A1-A2)
\end{rSection}

%----------------------------------------------------------------------------------------
%	INTERESTS SECTION
%----------------------------------------------------------------------------------------

\begin{rSection}{Interessen}
{\bf Fachlich}: Compilerbau, Maschinelle Sprachverarbeitung, Simulation, 3D-Rendering, KI \\
{\bf Persönlich}: Analoge Fotografie, Sprachen, Progressive Rock
\end{rSection}

\end{document}
